\documentclass{zhvt-classic}

% 使用本地源流明体（GenRyuMin2 TW）按文件路径加载，不依赖系统字体名
% 版框（内框／界栏／书口侧栏／鱼尾）改用与版心印章同一印泥棕色：
%   印章图片 yshs.png 实色 = RGB(135,68,52) = HTML 874434
\usepackage{xcolor}
\usepackage{tikz}
% sealbrown 双注册：xcolor 供 TikZ/pgf，L3 供 cls 版框
\definecolor{sealbrown}{HTML}{874434}
\ExplSyntaxOn
\color_set:nnn { sealbrown } { HTML } { 874434 }
\ExplSyntaxOff

\zhvtset{
  grid color       = sealbrown,
  main font        = GenRyuMin2TW-SB,
  main font path   = C:/Users/Administrator/texmf/fonts/opentype/public/GenRyuMin2TW/,
jiazhu font      = GenRyuMin2TW-M,
  jiazhu font path = C:/Users/Administrator/texmf/fonts/opentype/public/GenRyuMin2TW/,
  % 页码框容纳 4 字：正常页「九十七」+前置空 1 格=4 行恰好；
  % 100+ 页「一百零一」共 5 行超出，末字「一/二/三」被框裁掉（2026-10-04 实测），故放宽到 5
  shukou page number chars = 5,
  % 书口书名框高从默认 6 字减到 4 字，把省出的 2 字高拨给 chapter 框，
  % 使 6 字书名+卷名放进 9 字高的 chapter 框后不再压线（2026-10-05）
  shukou title chars = 4,
}

\title{史记}[匈奴列传]
\author{汉·司马迁}
\maker{闲叟漫记·重製}


% ===== 索隐述赞样式（2026-11-08 四次定稿 · 愚大人亲手改定参数）=====
% ⚠ 竖排前提：自定义字体家族必须带 RawFeature={vertical:+vert}；
%   simhei.ttf 加 +vhal 会整体加载失败静默回退，只写 +vert；simfang 同理。
\newCJKfontfamily{\shuzanfont}[Path=C:/Windows/Fonts/,RawFeature={vertical:+vert}]{simhei.ttf}
\newCJKfontfamily{\shuzanbody}[Path=C:/Windows/Fonts/,RawFeature={vertical:+vert}]{simfang.ttf}
% 题头：152x30pt 圆角 6pt 棕底白字；baseline (0,2) 使整框落入界行格子不压版框线；
% 白字节点 x 不补偿、y +8.2pt（像素级墨迹实测 dx/dy≈0，L/R 留白约 20pt）。
\newcommand{\shuzanhead}[1]{%
  \begin{tikzpicture}[baseline={(0,2)}]
    \fill[sealbrown, rounded corners=6pt]
      (-76pt,-15pt) rectangle (76pt,15pt);
    \node[text=white, inner sep=0pt, anchor=center] at (0pt,8.2pt)
      {\shuzanfont\fontsize{23}{23}\selectfont #1};
  \end{tikzpicture}%
}
% ★ 修正书口卷名位置：chapter 框起始位置改回 cls 默认值，避免全局影响。
%   只在每个卷首页通过 \juan 命令里的空命令微调该页位置（安全，不会跨页污染）。
\ExplSyntaxOn
\AtBeginDocument{
  \dim_set:Nn \l__zhvt_shukou_chapter_title_start_dim
    { \l__zhvt_shukou_title_height_dim + \l__zhvt_font_size_dim * 2 }
  \dim_set:Nn \l__zhvt_shukou_page_number_start_dim
    { \l__zhvt_shukou_chapter_title_start_dim + \l__zhvt_shukou_chapter_title_height_dim }
}
\ExplSyntaxOff

\begin{document}

% ★★ 版心（书口）排布覆写：把书名放到鱼尾之下（2026-10-04）
%
%   cls 默认版心自上而下：书名（title 静态框，在鱼尾之上）→ 鱼尾 → 篇名（chapter 动态框）→ 页码
%   本模板改为：          鱼尾（位置不动）→ 书名 → 篇名 → 页码分割线 → 页码
%
%   手法：把 cls 放在 title 静态框里的书名，改放进 chapter 动态框（鱼尾之下）；
%         title 静态框留空。书名取 \l__zhvt_title_str，即 \title{...} 的第一参数，
%         所以换书时只改 \title，不必在此重复写书名。
%
%   ⚠ 三条必须遵守的写法（2026-10-04 实测）：
%     ① 汉字必须用 \hbox{...} 包住。flowfram 会把框内容装进一个宽 14.4mm（= 一个 line_sep
%        = 书口栏宽）的纵向 minipage 并按该宽度折行，而整页靠 PDF /Rotate 90 显示、
%        框内文字本身不变向 —— 不包 \hbox 时「围炉夜话」4 字会被逐字折成 4 行，
%        在版面上就成了横排 4 字并横向溢出版框。
%     ② 竖排下**一个 \hbox 就是一整列**。书名与篇名要同处一列，必须放进同一个 \hbox；
%        分成两个 \hbox 时，后一个会被当成新的一列、甩到隔壁栏去
%        （实测「序」跑到书口栏外侧 14.4mm 的正文列上）。中间的 \quad 在框内是横向空白，
%        显示出来正好是书名与篇名之间的纵向间隔。
%     ③ ★ 内容盒要写成 **与框等宽的定宽盒** `\hbox to \l__zhvt_shukou_chapter_title_height_dim{…\hfil}`。
%        实测：chapter 框是把内容盒沿纵向**居中**摆放的，内容宽度一变，书名就跟着漂：
%        自然宽度的 \hbox（书名 38.4mm）落在 x=157.2mm，加篇名后（57.6mm）落在 x=152.4mm，
%        目录页（无篇名）与正文页的书名就差半格。做成等宽盒后，无论有没有篇名，
%        内容固定从框顶起排，实测两页均落在 x=142.8mm，不再漂。
\ExplSyntaxOn
% 把书口 chapter 框（书名+篇名）起始位置下移 1 字高，避免 4 字卷名顶到鱼尾
% ⚠ cls 的 \__zhvt_calculate_dimensions: 在 \AtBeginDocument 里执行，会重算该 dim；
%   必须把覆写挂到 begin document 之后，否则会被二次覆盖（2026-10-05 实测）。
% (旧方案已移到导言区，直接修改 chapter 框起始位置)
% 卷标题命令 #1=书口篇名(鱼尾) #2=正文卷名（含卷序，卷名内空格用 ~）
%
% ★★ 目录页两项修订（2026-10-05，愚大人指出）★★
%   ① 页码起点统一对齐：
%      原状 —— 目录项写成 \par\hspace{4em}卷名\hspace{4em}页码，页码起点 = 卷名实宽 + 4em。
%      卷名实宽随字数变：「卷一~梁惠王上」6 字 +~、「卷七~离娄上」5 字 +~，
%      差 1 个字位 → 卷七~卷十 4 条的页码整体比其余 10 条高 27.21pt（一个字），看着上下不齐。
%      改法 —— 卷名补全角空格（宽 = 1 字 = 27.215pt）补足到 **7 字符位**再写进目录，
%      14 条页码起点遂全部落在同一条竖线上（起点对齐；末字随页码字数自然参差）。
%      ⚠ 只对目录里的副本 \l_tmpc_tl 补，正文卷名仍用 \l_tmpa_tl 原样 —— 正文像素不受影响。
%   ② 卷名整体上移半字位：前导 \hspace 由 4em 改 3.5em。
%      实测原状目录块在版框内偏下（天头留白 139.6pt / 地脚留白 110.7pt），
%      上移 0.5em = 13.6pt 后上下基本均分，视觉居中。
\NewDocumentCommand { \juan } { mm } {
	\clearpage
	\tl_set:Nn \l__zhvt_chapter_title_tl { #1 }
	\tl_set:Nn \l_tmpa_tl { #2 }
	\tl_set_eq:NN \l_tmpc_tl \l_tmpa_tl
	\int_step_inline:nn { 7 - \str_count:n { #2 } }
	  { \tl_put_right:Nn \l_tmpc_tl { 　 } }
	\tl_set:Nx \l_tmpb_tl { \zhnum { page } }
	\tl_put_left:Nn \l_tmpb_tl { \hspace { 4em } }
	\tl_put_left:Nn \l_tmpb_tl { \l_tmpc_tl }
	\tl_put_left:Nn \l_tmpb_tl { \par \hspace { 3.5em } }
	\addtocontents { toc } { \l_tmpb_tl }
	\skip_horizontal:n { \l__zhvt_font_size_dim * 4 }
	\l_tmpa_tl
	\par
	\thispagestyle { empty }
}
\cs_set_protected:Npn \__zhvt_set_shukou_content: {
	% ★ 鱼尾上方留白（title 静态框）嵌入闲叟钤印：
	%   整页 /Rotate 90，故此处把图用 \rotatebox{90} 转正为竖章；
	%   尺寸卡在版心条幅内（书内约 7mm 宽 × 33.8mm 高，留余量不压鱼尾）。
	\setstaticcontents * { title } {
		\vbox_to_ht:nn { \l__zhvt_baseline_skip_dim } {
			\vskip 14.2pt
			\hbox_to_wd:nn { \l__zhvt_shukou_title_height_dim } {
				\hfil
				\rotatebox [ origin = c ] { 90 } { \includegraphics [ width = 7mm , height = 33.8mm ] { yshs.png } }
				\hfil
			}
		}
	}
% ★ 版心（书口）统一改法（2026-11-16 愚大人定）：18pt＋框内居中＋书口文字同版框色
%   ① 原 27.2pt 时 7 字长篇名（如「廉颇蔺相如列传」）末字溢出 chapter 框、压住页码；
%      改 \fontsize{18pt}{18pt}，前后 \hfil 令字串在框内居中（不随篇名字数漂移）。
%   ② \vbox{\vskip 8.2pt} 定量下移，抵消缩字号后栏宽方向 8.2pt 上偏（像素实测 y 中心 543.3 = 原版）。
%      ⚠ \raisebox / \vspace* 在 flowfram 框内无效；换字号须同步换补偿：16pt→10.1、19pt→7.4。
%   ③ chapter 与 page~number 均置 \color{sealbrown}（HTML 874434），与版框／界行同色；正文仍黑。
%   ⚠ 页码字位、钤印、鱼尾、界行、正文一律未动；五部合订本同此口径。
	\setdynamiccontents * { chapter }
	{ \color{sealbrown} \vbox { \vskip 8.2pt \hbox to \l__zhvt_shukou_chapter_title_height_dim
		{ \hfil \fontsize{18pt}{18pt}\selectfont \l__zhvt_title_str \quad \l__zhvt_chapter_title_tl \hfil } } }
	\setdynamiccontents * { page~number } { \color{sealbrown} \quad \zhnum { page } }
}
\__zhvt_set_shukou_content:
\ExplSyntaxOff

\maketitle

% 单篇，无多卷目录，不渲染目录页

\juan{匈奴列传}{卷一百十~匈奴列传第五十}
\hspace*{2em}匈奴。其先祖夏后氏之苗裔也。曰淳维。唐虞以上有山戎。獫狁。荤粥。居于北蛮。随畜牧而转移。其畜之所多则马。牛。羊。其奇畜则橐扆。驴骡。駃騠。騊駼。驒騱。逐水草迁徙。毋城郭常处耕田之业。然亦各有分地。毋文书。以言语为约束。兒能骑羊。引弓射鸟鼠。少长则射狐兔。用为食。士力能毌弓。尽为甲骑。其俗。宽则随畜。因射猎禽兽为生业。急则人习战攻以侵伐。其天性也。其长兵则弓矢。短兵则刀鋋。利则进。不利则退。不羞遁走。苟利所在。不知礼义。自君王以下。咸食畜肉。衣其皮革。被旃裘。壮者食肥美。老者食其余。贵壮健。贱老弱。父死。妻其后母。兄弟死。皆取其妻妻之。其俗有名不讳。而无姓字。\\
\hspace*{2em}夏道衰。而公刘失其稷官。变于西戎。邑于豳。其后三百有余岁。戎狄攻大王亶父。亶父亡走岐下。而豳人悉从亶父而邑焉。作周。其后百有余岁。周西伯昌伐畎夷氏。后十有余年。武王伐纣而营雒邑。复居于酆鄗。放逐戎夷泾。洛之北。以时入贡。命曰﹁荒服﹂。其后二百有余年。周道衰。而穆王伐犬戎。得四白狼四白鹿以归。自是之后。荒服不至。於是周遂作甫刑之辟。穆王之后二百有余年。周幽王用宠姬襃姒之故。与申侯有卻。申侯怒而与犬戎共攻杀周幽王于骊山之下。遂取周之焦穫。而居于泾渭之间。侵暴中国。秦襄公救周。於是周平王去酆鄗而东徙雒邑。当是之时。秦襄公伐戎至岐。始列为诸侯。是后六十有五年。而山戎越燕而伐齐。齐釐公与战于齐郊。其后四十四年。而山戎伐燕。燕告急于齐。齐桓公北伐山戎。山戎走。其后二十有余年。而戎狄至洛邑。伐周襄王。襄王奔于郑之氾邑。初。周襄王欲伐郑。故娶戎狄女为后。与戎狄兵共伐郑。已而黜狄后。狄后怨。而襄王后母曰惠后。有子子带。欲立之。於是惠后与狄后。子带为内应。开戎狄。戎狄以故得入。破逐周襄王。而立子带为天子。於是戎狄或居于陆浑。东至於卫。侵盗暴虐中国。中国疾之。故诗人歌之曰﹁戎狄是应﹂。﹁薄伐獫狁。至於大原﹂。﹁出舆彭彭。城彼朔方﹂。周襄王既居外四年。乃使使告急于晋。晋文公初立。欲修霸业。乃兴师伐逐戎翟。诛子带。迎内周襄王。居于雒邑。\\
\hspace*{2em}当是之时。秦晋为彊国。晋文公攘戎翟。居于河西。洛之间。号曰赤翟。白翟。秦穆公得由余。西戎八国服於秦。故自陇以西有绵诸。绲戎。翟。镕之戎。岐。梁山。泾。漆之北有义渠。大荔。乌氏。朐衍之戎。而晋北有林胡。楼烦之戎。燕北有东胡。山戎。各分散居谿谷。自有君长。往往而聚者百有余戎。然莫能相一。\\
\hspace*{2em}自是之后百有余年。晋悼公使魏绛和戎翟。戎翟朝晋。后百有余年。赵襄子逾句注而破并代以临胡貉。其后既与韩魏共灭智伯。分晋地而有之。则赵有代。句注之北。魏有河西。上郡。以与戎界边。其后义渠之戎筑城郭以自守。而秦稍蚕食。至於惠王。遂拔义渠二十五城。惠王击魏。魏尽入西河及上郡于秦。秦昭王时。义渠戎王与宣太后乱。有二子。宣太后诈而杀义渠戎王於甘泉。遂起兵伐残义渠。於是秦有陇西。北地。上郡。筑长城以拒胡。而赵武灵王亦变俗胡服。习骑射。北破林胡。楼烦。筑长城。自代并阴山下。至高阙为塞。而置云中。雁门。代郡。其后燕有贤将秦开。为质於胡。胡甚信之。归而袭破走东胡。东胡卻千余里。与荆轲刺秦王秦舞阳者。开之孙也。燕亦筑长城。自造阳至襄平。置上谷。渔阳。右北平。辽西。辽东郡以拒胡。当是之时。冠带战国七。而三国边於匈奴。其后赵将李牧时。匈奴不敢入赵边。后秦灭六国。而始皇帝使蒙恬将十万之众北击胡。悉收河南地。因河为塞。筑四十四县城临河。徙適戍以充之。而通直道。自九原至云阳。因边山险巉谿谷可缮者治之。起临洮至辽东万余里。又度河据阳山北假中。\\
\hspace*{2em}当是之时。东胡彊而月氏盛。匈奴单于曰头曼。头曼不胜秦。北徙。十余年而蒙恬死。诸侯畔秦。中国扰乱。诸秦所徙適戍边者皆复去。於是匈奴得宽。复稍度河南与中国界於故塞。\\
\hspace*{2em}单于有太子名冒顿。后有所爱阏氏。生少子。而单于欲废冒顿而立少子。乃使冒顿质於月氏。冒顿既质於月氏。而头曼急击月氏。月氏欲杀冒顿。冒顿盗其善马。骑之亡归。头曼以为壮。令将万骑。冒顿乃作为鸣镝。习勒其骑射。令曰。﹁鸣镝所射而不悉射者。斩之。﹂行猎鸟兽。有不射鸣镝所射者。辄斩之。已而冒顿以鸣镝自射其善马。左右或不敢射者。冒顿立斩不射善马者。居顷之。复以鸣镝自射其爱妻。左右或颇恐。不敢射。冒顿又复斩之。居顷之。冒顿出猎。以鸣镝射单于善马。左右皆射之。於是冒顿知其左右皆可用。从其父单于头曼猎。以鸣镝射头曼。其左右亦皆随鸣镝而射杀单于头曼。遂尽诛其后母与弟及大臣不听从者。冒顿自立为单于。\\
\hspace*{2em}冒顿既立。是时东胡彊盛。闻冒顿杀父自立。乃使使谓冒顿。欲得头曼时有千里马。冒顿问群臣。群臣皆曰。﹁千里马。匈奴宝马也。勿与。﹂冒顿曰。﹁柰何与人邻国而爱一马乎。﹂遂与之千里马。居顷之。东胡以为冒顿畏之。乃使使谓冒顿。欲得单于一阏氏。冒顿复问左右。左右皆怒曰。﹁东胡无道。乃求阏氏。请击之。﹂冒顿曰。﹁柰何与人邻国爱一女子乎。﹂遂取所爱阏氏予东胡。东胡王愈益骄。西侵。与匈奴间。中有弃地。莫居。千余里。各居其边为瓯脱。东胡使使谓冒顿曰。﹁匈奴所与我界瓯脱外弃地。匈奴非能至也。吾欲有之。﹂冒顿问群臣。群臣或曰。﹁此弃地。予之亦可。勿予亦可。﹂於是冒顿大怒曰。﹁地者。国之本也。柰何予之。﹂诸言予之者。皆斩之。冒顿上马。令国中有后者斩。遂东袭击东胡。东胡初轻冒顿。不为备。及冒顿以兵至。击。大破灭东胡王。而虏其民人及畜产。既归。西击走月氏。南并楼烦。白羊河南王。悉复收秦所使蒙恬所夺匈奴地者。与汉关故河南塞。至朝那。肤施。遂侵燕。代。是时汉兵与项羽相距。中国罢於兵革。以故冒顿得自彊。控弦之士三十余万。\\
\hspace*{2em}自淳维以至头曼千有余岁。时大时小。别散分离。尚矣。其世传不可得而次云。然至冒顿而匈奴最彊大。尽服从北夷。而南与中国为敌国。其世传国官号乃可得而记云。\\
\hspace*{2em}置左右贤王。左右谷蠡王。左右大将。左右大都尉。左右大当户。左右骨都侯。匈奴谓贤曰﹁屠耆﹂。故常以太子为左屠耆王。自如左右贤王以下至当户。大者万骑。小者数千。凡二十四长。立号曰﹁万骑﹂。诸大臣皆世官。呼衍氏。兰氏。其后有须卜氏。此三姓其贵种也。诸左方王将居东方。直上谷以往者。东接秽貉。朝鲜。右方王将居西方。直上郡以西。接月氏。氐。羌。而单于之庭直代。云中。各有分地。逐水草移徙。而左右贤王。左右谷蠡王最为大。左右骨都侯辅政。诸二十四长亦各自置千长。百长。什长。裨小王。相。封都尉。当户。且渠之属。\\
\hspace*{2em}岁正月。诸长小会单于庭。祠。五月。大会茏城。祭其先。天地。鬼神。秋。马肥。大会蹛林。课校人畜计。其法。拔刃尺者死。坐盗者没入其家。有罪小者轧。大者死。狱久者不过十日。一国之囚不过数人。而单于朝出营。拜日之始生。夕拜月。其坐。长左而北乡。日上戊己。其送死。有棺椁金银衣裘。而无封树丧服。近幸臣妾从死者。多至数千百人。举事而候星月。月盛壮则攻战。月亏则退兵。其攻战。斩首虏赐一卮酒。而所得卤获因以予之。得人以为奴婢。故其战。人人自为趣利。善为诱兵以冒敌。故其见敌则逐利。如鸟之集。其困败。则瓦解云散矣。战而扶舆死者。尽得死者家财。\\
\hspace*{2em}后北服浑庾。屈射。丁零。鬲昆。薪犁之国。於是匈奴贵人大臣皆服。以冒顿单于为贤。\\
\hspace*{2em}是时汉初定中国。徙韩王信於代。都马邑。匈奴大攻围马邑。韩王信降匈奴。匈奴得信。因引兵南逾句注。攻太原。至晋阳下。高帝自将兵往击之。会冬大寒雨雪。卒之堕指者十二三。於是冒顿详败走。诱汉兵。汉兵逐击冒顿。冒顿匿其精兵。见其羸弱。於是汉悉兵。多步兵。三十二万。北逐之。高帝先至平城。步兵未尽到。冒顿纵精兵四十万骑围高帝於白登。七日。汉兵中外不得相救饷。匈奴骑。其西方尽白马。东方尽青駹马。北方尽乌骊马。南方尽骍马。高帝乃使使间厚遗阏氏。阏氏乃谓冒顿曰。﹁两主不相困。今得汉地。而单于终非能居之也。且汉王亦有神。单于察之。﹂冒顿与韩王信之将王黄。赵利期。而黄。利兵又不来。疑其与汉有谋。亦取阏氏之言。乃解围之一角。於是高帝令士皆持满傅矢外乡。从解角直出。竟与大军合。而冒顿遂引兵而去。汉亦引兵而罢。使刘敬结和亲之约。\\
\hspace*{2em}是后韩王信为匈奴将。及赵利。王黄等数倍约。侵盗代。云中。居无几何。陈豨反。又与韩信合谋击代。汉使樊哙往击之。复拔代。雁门。云中郡县。不出塞。是时匈奴以汉将众往降。故冒顿常往来侵盗代地。於是汉患之。高帝乃使刘敬奉宗室女公主为单于阏氏。岁奉匈奴絮缯酒米食物各有数。约为昆弟以和亲。冒顿乃少止。后燕王卢绾反。率其党数千人降匈奴。往来苦上谷以东。\\
\hspace*{2em}高祖崩。孝惠。吕太后时。汉初定。故匈奴以骄。冒顿乃为书遗高后。妄言。高后欲击之。诸将曰。﹁以高帝贤武。然尚困於平城。﹂於是高后乃止。复与匈奴和亲。\\
\hspace*{2em}至孝文帝初立。复修和亲之事。其三年五月。匈奴右贤王入居河南地。侵盗上郡葆塞蛮夷。杀略人民。於是孝文帝诏丞相灌婴发车骑八万五千。诣高奴。击右贤王。右贤王走出塞。文帝幸太原。是时济北王反。文帝归。罢丞相击胡之兵。\\
\hspace*{2em}其明年。单于遗汉书曰。﹁天所立匈奴大单于敬问皇帝无恙。前时皇帝言和亲事。称书意。合欢。汉边吏侵侮右贤王。右贤王不请。听后义卢侯难氏等计。与汉吏相距。绝二主之约。离兄弟之亲。皇帝让书再至。发使以书报。不来。汉使不至。汉以其故不和。邻国不附。今以小吏之败约故。罚右贤王。使之西求月氏击之。以天之福。吏卒良。马彊力。以夷灭月氏。尽斩杀降下之。定楼兰。乌孙。呼揭及其旁二十六国。皆以为匈奴。诸引弓之民。并为一家。北州已定。原寝兵休士卒养马。除前事。复故约。以安边民。以应始古。使少者得成其长。老者安其处。世世平乐。未得皇帝之志也。故使郎中系雩浅奉书请。献橐他一匹。骑马二匹。驾二驷。皇帝即不欲匈奴近塞。则且诏吏民远舍。使者至。即遣之。﹂以六月中来至薪望之地。书至。汉议击与和亲孰便。公卿皆曰。﹁单于新破月氏。乘胜。不可击。且得匈奴地。泽卤。非可居也。和亲甚便。﹂汉许之。\\
\hspace*{2em}孝文皇帝前六年。汉遗匈奴书曰。﹁皇帝敬问匈奴大单于无恙。使郎中系雩浅遗朕书曰。﹃右贤王不请。听后义卢侯难氏等计。绝二主之约。离兄弟之亲。汉以故不和。邻国不附。今以小吏败约。故罚右贤王使西击月氏。尽定之。原寝兵休士卒养马。除前事。复故约。以安边民。使少者得成其长。老者安其处。世世平乐。﹄朕甚嘉之。此古圣主之意也。汉与匈奴约为兄弟。所以遗单于甚厚。倍约离兄弟之亲者。常在匈奴。然右贤王事已在赦前。单于勿深诛。单于若称书意。明告诸吏。使无负约。有信。敬如单于书。使者言单于自将伐国有功。甚苦兵事。服绣袷绮衣。绣袷长襦。锦袷袍各一。比余一。黄金饰具带一。黄金胥纰一。绣十匹。锦三十匹。赤綈。绿缯各四十匹。使中大夫意。谒者令肩遗单于。﹂\\
\hspace*{2em}后顷之。冒顿死。子稽粥立。号曰老上单于。\\
\hspace*{2em}老上稽粥单于初立。孝文皇帝复遣宗室女公主为单于阏氏。使宦者燕人中行说傅公主。说不欲行。汉彊使之。说曰。﹁必我行也。为汉患者。﹂中行说既至。因降单于。单于甚亲幸之。\\
\hspace*{2em}初。匈奴好汉缯絮食物。中行说曰。﹁匈奴人众不能当汉之一郡。然所以彊者。以衣食异。无仰於汉也。今单于变俗好汉物。汉物不过什二。则匈奴尽归於汉矣。其得汉缯絮。以驰草棘中。衣袴皆裂敝。以示不如旃裘之完善也。得汉食物皆去之。以示不如湩酪之便美也。﹂於是说教单于左右疏记。以计课其人众畜物。\\
\hspace*{2em}汉遗单于书。牍以尺一寸。辞曰﹁皇帝敬问匈奴大单于无恙﹂。所遗物及言语云云。中行说令单于遗汉书以尺二寸牍。及印封皆令广大长。倨傲其辞曰﹁天地所生日月所置匈奴大单于敬问汉皇帝无恙﹂。所以遗物言语亦云云。\\
\hspace*{2em}汉使或言曰。﹁匈奴俗贱老。﹂中行说穷汉使曰。﹁而汉俗屯戍从军当发者。其老亲岂有不自脱温厚肥美以赍送饮食行戍乎。﹂汉使曰。﹁然。﹂中行说曰。﹁匈奴明以战攻为事。其老弱不能斗。故以其肥美饮食壮健者。盖以自为守卫。如此父子各得久相保。何以言匈奴轻老也。﹂汉使曰。﹁匈奴父子乃同穹庐而卧。父死。妻其后母。兄弟死。尽取其妻妻之。无冠带之饰。阙庭之礼。﹂中行说曰。﹁匈奴之俗。人食畜肉。饮其汁。衣其皮。畜食草饮水。随时转移。故其急则人习骑射。宽则人乐无事。其约束轻。易行也。君臣简易。一国之政犹一身也。父子兄弟死。取其妻妻之。恶种姓之失也。故匈奴虽乱。必立宗种。今中国虽详不取其父兄之妻。亲属益疏则相杀。至乃易姓。皆从此类。且礼义之敝。上下交怨望。而室屋之极。生力必屈。夫力耕桑以求衣食。筑城郭以自备。故其民急则不习战功。缓则罢於作业。嗟土室之人。顾无多辞。令喋喋而佔々。冠固何当。﹂\\
\hspace*{2em}自是之后。汉使欲辩论者。中行说辄曰。﹁汉使无多言。顾汉所输匈奴缯絮米糵。令其量中。必善美而己矣。何以为言乎。且所给备善则已。不备。苦恶。则候秋孰。以骑驰蹂而稼穑耳。﹂日夜教单于候利害处。\\
\hspace*{2em}汉孝文皇帝十四年。匈奴单于十四万骑入朝那。萧关。杀北地都尉卬。虏人民畜产甚多。遂至彭阳。使奇兵入烧回中宫。候骑至雍甘泉。於是文帝以中尉周舍。郎中令张武为将军。发车千乘。骑十万。军长安旁以备胡寇。而拜昌侯卢卿为上郡将军。甯侯魏遫为北地将军。隆虑侯周灶为陇西将军。东阳侯张相如为大将军。成侯董赤为前将军。大发车骑往击胡。单于留塞内月余乃去。汉逐出塞即还。不能有所杀。匈奴日已骄。岁入边。杀略人民畜产甚多。云中。辽东最甚。至代郡万余人。汉患之。乃使使遗匈奴书。单于亦使当户报谢。复言和亲事。\\
\hspace*{2em}孝文帝后二年。使使遗匈奴书曰。﹁皇帝敬问匈奴大单于无恙。使当户且居雕渠难。郎中韩辽遗朕马二匹。已至。敬受。先帝制。长城以北。引弓之国。受命单于。长城以内。冠带之室。朕亦制之。使万民耕织射猎衣食。父子无离。臣主相安。俱无暴逆。今闻渫恶民贪降其进取之利。倍义绝约。忘万民之命。离两主之驩。然其事已在前矣。书曰。﹃二国已和亲。两主驩说。寝兵休卒养马。世世昌乐。闟然更始。﹄朕甚嘉之。圣人者日新。改作更始。使老者得息。幼者得长。各保其首领而终其天年。朕与单于俱由此道。顺天恤民。世世相传。施之无穷。天下莫不咸便。汉与匈奴邻国之敌。匈奴处北地。寒。杀气早降。故诏吏遗单于秫糵金帛丝絮佗物岁有数。今天下大安。万民熙熙。朕与单于为之父母。朕追念前事。薄物细故。谋臣计失。皆不足以离兄弟之驩。朕闻天不颇覆。地不偏载。朕与单于皆捐往细故。俱蹈大道。堕坏前恶。以图长久。使两国之民若一家子。元元万民。下及鱼鳖。上及飞鸟。跂行喙息蠕动之类。莫不就安利而辟危殆。故来者不止。天之道也。俱去前事。朕释逃虏民。单于无言章尼等。朕闻古之帝王。约分明而无食言。单于留志。天下大安。和亲之后。汉过不先。单于其察之。﹂\\
\hspace*{2em}单于既约和亲。於是制诏御史曰。﹁匈奴大单于遗朕书。言和亲已定。亡人不足以益众广地。匈奴无入塞。汉无出塞。犯约者杀之。可以久亲。后无咎。俱便。朕已许之。其布告天下。使明知之。﹂\\
\hspace*{2em}后四岁。老上稽粥单于死。子军臣立为单于。既立。孝文皇帝复与匈奴和亲。而中行说复事之。\\
\hspace*{2em}军臣单于立四岁。匈奴复绝和亲。大入上郡。云中各三万骑。所杀略甚众而去。於是汉使三将军军屯北地。代屯句注。赵屯飞狐口。缘边亦各坚守以备胡寇。又置三将军。军长安西细柳。渭北棘门。霸上以备胡。胡骑入代句注边。烽火通於甘泉。长安。数月。汉兵至边。匈奴亦去远塞。汉兵亦罢。后岁余。孝文帝崩。孝景帝立。而赵王遂乃阴使人於匈奴。吴楚反。欲与赵合谋入边。汉围破赵。匈奴亦止。自是之后。孝景帝复与匈奴和亲。通关市。给遗匈奴。遣公主。如故约。终孝景时。时小入盗边。无大寇。\\
\hspace*{2em}今帝即位。明和亲约束。厚遇。通关市。饶给之。匈奴自单于以下皆亲汉。往来长城下。\\
\hspace*{2em}汉使马邑下人聂翁壹奸兰出物与匈奴交。详为卖马邑城以诱单于。单于信之。而贪马邑财物。乃以十万骑入武州塞。汉伏兵三十余万马邑旁。御史大夫韩安国为护军。护四将军以伏单于。单于既入汉塞。未至马邑百余里。见畜布野而无人牧者。怪之。乃攻亭。是时雁门尉史行徼。见寇。葆此亭。知汉兵谋。单于得。欲杀之。尉史乃告单于汉兵所居。单于大惊曰。﹁吾固疑之。﹂乃引兵还。出曰。﹁吾得尉史。天也。天使若言。﹂以尉史为﹁天王﹂。汉兵约单于入马邑而纵。单于不至。以故汉兵无所得。汉将军王恢部出代击胡辎重。闻单于还。兵多。不敢出。汉以恢本造兵谋而不进。斩恢。自是之后。匈奴绝和亲。攻当路塞。往往入盗於汉边。不可胜数。然匈奴贪。尚乐关市。嗜汉财物。汉亦尚关市不绝以中之。\\
\hspace*{2em}自马邑军后五年之秋。汉使四将军各万骑击胡关市下。将军卫青出上谷。至茏城。得胡首虏七百人。公孙贺出云中。无所得。公孙敖出代郡。为胡所败七千余人。李广出雁门。为胡所败。而匈奴生得广。广后得亡归。汉囚敖。广。敖。广赎为庶人。其冬。匈奴数入盗边。渔阳尤甚。汉使将军韩安国屯渔阳备胡。其明年秋。匈奴二万骑入汉。杀辽西太守。略二千余人。胡又入败渔阳太守军千余人。围汉将军安国。安国时千余骑亦且尽。会燕救至。匈奴乃去。匈奴又入雁门。杀略千余人。於是汉使将军卫青将三万骑出雁门。李息出代郡。击胡。得首虏数千人。其明年。卫青复出云中以西至陇西。击胡之楼烦。白羊王於河南。得胡首虏数千。牛羊百余万。於是汉遂取河南地。筑朔方。复缮故秦时蒙恬所为塞。因河为固。汉亦弃上谷之什辟县造阳地以予胡。是岁。汉之元朔二年也。\\
\hspace*{2em}其后冬。匈奴军臣单于死。军臣单于弟左谷蠡王伊稚斜自立为单于。攻破军臣单于太子於单。於单亡降汉。汉封於单为涉安侯。数月而死。\\
\hspace*{2em}伊稚斜单于既立。其夏。匈奴数万骑入杀代郡太守恭友。略千余人。其秋。匈奴又入雁门。杀略千余人。其明年。匈奴又复复入代郡。定襄。上郡。各三万骑。杀略数千人。匈奴右贤王怨汉夺之河南地而筑朔方。数为寇。盗边。及入河南。侵扰朔方。杀略吏民其众。\\
\hspace*{2em}其明年春。汉以卫青为大将军。将六将军。十余万人。出朔方。高阙击胡。右贤王以为汉兵不能至。饮酒醉。汉兵出塞六七百里。夜围右贤王。右贤王大惊。脱身逃走。诸精骑往往随后去。汉得右贤王众男女万五千人。裨小王十余人。其秋。匈奴万骑入杀代郡都尉硃英。略千余人。\\
\hspace*{2em}其明年春。汉复遣大将军卫青将六将军。兵十余万骑。乃再出定襄数百里击匈奴。得首虏前后凡万九千余级。而汉亦亡两将军。军三千余骑。右将军建得以身脱。而前将军翕侯赵信兵不利。降匈奴。赵信者。故胡小王。降汉。汉封为翕侯。以前将军与右将军并军分行。独遇单于兵。故尽没。单于既得翕侯。以为自次王。用其姊妻之。与谋汉。信教单于益北绝幕。以诱罢汉兵。徼极而取之。无近塞。单于从其计。其明年。胡骑万人入上谷。杀数百人。\\
\hspace*{2em}其明年春。汉使骠骑将军去病将万骑出陇西。过焉支山千余里。击匈奴。得胡首虏万八千余级。破得休屠王祭天金人。其夏。骠骑将军复与合骑侯数万骑出陇西。北地二千里。击匈奴。过居延。攻祁连山。得胡首虏三万余人。裨小王以下七十余人。是时匈奴亦来入代郡。雁门。杀略数百人。汉使博望侯及李将军广出右北平。击匈奴左贤王。左贤王围李将军。卒可四千人。且尽。杀虏亦过当。会博望侯军救至。李将军得脱。汉失亡数千人。合骑侯后骠骑将军期。及与博望侯皆当死。赎为庶人。\\
\hspace*{2em}其秋。单于怒浑邪王。休屠王居西方为汉所杀虏数万人。欲召诛之。浑邪王与休屠王恐。谋降汉。汉使骠骑将军往迎之。浑邪王杀休屠王。并将其众降汉。凡四万余人。号十万。於是汉已得浑邪王。则陇西。北地。河西益少胡寇。徙关东贫民处所夺匈奴河南。新秦中以实之。而减北地以西戍卒半。其明年。匈奴入右北平。定襄各数万骑。杀略千余人而去。\\
\hspace*{2em}其明年春。汉谋曰﹁翕侯信为单于计。居幕北。以为汉兵不能至﹂。乃粟马发十万骑。私从马凡十四万匹。粮重不与焉。令大将军青。骠骑将军去病中分军。大将军出定襄。骠骑将军出代。咸约绝幕击匈奴。单于闻之。远其辎重。以精兵待於幕北。与汉大将军接战一日。会暮。大风起。汉兵纵左右翼围单于。单于自度战不能如汉兵。单于遂独身与壮骑数百溃汉围西北遁走。汉兵夜追不得。行斩捕匈奴首虏万九千级。北至阗颜山赵信城而还。\\
\hspace*{2em}单于之遁走。其兵往往与汉兵相乱而随单于。单于久不与其大众相得。其右谷蠡王以为单于死。乃自立为单于。真单于复得其众。而右谷蠡王乃去其单于号。复为右谷蠡王。\\
\hspace*{2em}汉骠骑将军之出代二千余里。与左贤王接战。汉兵得胡首虏凡七万余级。左贤王将皆遁走。骠骑封於狼居胥山。禅姑衍。临翰海而还。\\
\hspace*{2em}是后匈奴远遁。而幕南无王庭。汉度河自朔方以西至令居。往往通渠置田。官吏卒五六万人。稍蚕食。地接匈奴以北。\\
\hspace*{2em}初。汉两将军大出围单于。所杀虏八九万。而汉士卒物故亦数万。汉马死者十余万。匈奴虽病。远去。而汉亦马少。无以复往。匈奴用赵信之计。遣使於汉。好辞请和亲。天子下其议。或言和亲。或言遂臣之。丞相长史任敞曰。﹁匈奴新破。困。宜可使为外臣。朝请於边。﹂汉使任敞於单于。单于闻敞计。大怒。留之不遣。先是汉亦有所降匈奴使者。单于亦辄留汉使相当。汉方复收士马。会骠骑将军去病死。於是汉久不北击胡。\\
\hspace*{2em}数岁。伊稚斜单于立十三年死。子乌维立为单于。是岁。汉元鼎三年也。乌维单于立。而汉天子始出巡郡县。其后汉方南诛两越。不击匈奴。匈奴亦不侵入边。\\
\hspace*{2em}乌维单于立三年。汉已灭南越。遣故太仆贺将万五千骑出九原二千余里。至浮苴井而还。不见匈奴一人。汉又遣故从骠侯赵破奴万余骑出令居数千里。至匈河水而还。亦不见匈奴一人。\\
\hspace*{2em}是时天子巡边。至朔方。勒兵十八万骑以见武节。而使郭吉风告单于。郭吉既至匈奴。匈奴主客问所使。郭吉礼卑言好。曰。﹁吾见单于而口言。﹂单于见吉。吉曰。﹁南越王头已悬於汉北阙。今单于即前与汉战。天子自将兵待边。单于即不能。即南面而臣於汉。何徒远走。亡匿於幕北寒苦无水草之地。毋为也。﹂语卒而单于大怒。立斩主客见者。而留郭吉不归。迁之北海上。而单于终不肯为寇於汉边。休养息士马。习射猎。数使使於汉。好辞甘言求请和亲。\\
\hspace*{2em}汉使王乌等窥匈奴。匈奴法。汉使非去节而以墨黥其面者不得入穹庐。王乌。北地人。习胡俗。去其节。黥面。得入穹庐。单于爱之。详许甘言。为遣其太子入汉为质。以求和亲。\\
\hspace*{2em}汉使杨信於匈奴。是时汉东拔秽貉。朝鲜以为郡。而西置酒泉郡以鬲绝胡与羌通之路。汉又西通月氏。大夏。又以公主妻乌孙王。以分匈奴西方之援国。又北益广田至胘雷为塞。而匈奴终不敢以为言。是岁。翕侯信死。汉用事者以匈奴为已弱。可臣从也。杨信为人刚直屈彊。素非贵臣。单于不亲。单于欲召入。不肯去节。单于乃坐穹庐外见杨信。杨信既见单于。说曰。﹁即欲和亲。以单于太子为质於汉。﹂单于曰。﹁非故约。故约。汉常遣翁主。给缯絮食物有品。以和亲。而匈奴亦不扰边。今乃欲反古。令吾太子为质。无几矣。﹂匈奴俗。见汉使非中贵人。其儒先。以为欲说。折其辩。其少年。以为欲刺。折其气。每汉使入匈奴。匈奴辄报偿。汉留匈奴使。匈奴亦留汉使。必得当乃肯止。\\
\hspace*{2em}杨信既归。汉使王乌。而单于复谩以甘言。欲多得汉财物。绐谓王乌曰。﹁吾欲入汉见天子。面相约为兄弟。﹂王乌归报汉。汉为单于筑邸于长安。匈奴曰。﹁非得汉贵人使。吾不与诚语。﹂匈奴使其贵人至汉。病。汉予药。欲愈之。不幸而死。而汉使路充国佩二千石印绶往使。因送其丧。厚葬直数千金。曰﹁此汉贵人也﹂。单于以为汉杀吾贵使者。乃留路充国不归。诸所言者。单于特空绐王乌。殊无意入汉及遣太子来质。於是匈奴数使奇兵侵犯边。汉乃拜郭昌为拔胡将军。及浞野侯屯朔方以东。备胡。路充国留匈奴三岁。单于死。\\
\hspace*{2em}乌维单于立十岁而死。子乌师庐立为单于。年少。号为兒单于。是岁元封六年也。自此之后。单于益西北。左方兵直云中。右方直酒泉。燉煌郡。\\
\hspace*{2em}兒单于立。汉使两使者。一吊单于。一吊右贤王。欲以乖其国。使者入匈奴。匈奴悉将致单于。单于怒而尽留汉使。汉使留匈奴者前后十余辈。而匈奴使来。汉亦辄留相当。\\
\hspace*{2em}是岁。汉使贰师将军广利西伐大宛。而令因杅将军敖筑受降城。其冬。匈奴大雨雪。畜多饥寒死。兒单于年少。好杀伐。国人多不安。左大都尉欲杀单于。使人间告汉曰。﹁我欲杀单于降汉。汉远。即兵来迎我。我即发。﹂初。汉闻此言。故筑受降城。犹以为远。\\
\hspace*{2em}其明年春。汉使浞野侯破奴将二万余骑出朔方西北二千余里。期至浚稽山而还。浞野侯既至期而还。左大都尉欲发而觉。单于诛之。发左方兵击浞野。浞野侯行捕首虏得数千人。还。未至受降城四百里。匈奴兵八万骑围之。浞野侯夜自出求水。匈奴间捕。生得浞野侯。因急击其军。军中郭纵为护。维王为渠。相与谋曰。﹁及诸校尉畏亡将军而诛之。莫相劝归。﹂军遂没於匈奴。匈奴兒单于大喜。遂遣奇兵攻受降城。不能下。乃寇入边而去。其明年。单于欲自攻受降城。未至。病死。\\
\hspace*{2em}兒单于立三岁而死。子年少。匈奴乃立其季父乌维单于弟右贤王呴犁湖为单于。是岁太初三年也。\\
\hspace*{2em}呴犁湖单于立。汉使光禄徐自为出五原塞数百里。远者千余里。筑城鄣列亭至庐朐。而使游击将军韩说。长平侯卫伉屯其旁。使彊弩都尉路博德筑居延泽上。\\
\hspace*{2em}其秋。匈奴大入定襄。云中。杀略数千人。败数二千石而去。行破坏光禄所筑城列亭鄣。又使右贤王入酒泉。张掖。略数千人。会任文击救。尽复失所得而去。是岁。贰师将军破大宛。斩其王而还。匈奴欲遮之。不能至。其冬。欲攻受降城。会单于病死。\\
\hspace*{2em}呴犁湖单于立一岁死。匈奴乃立其弟左大都尉且鞮侯为单于。\\
\hspace*{2em}汉既诛大宛。威震外国。天子意欲遂困胡。乃下诏曰。﹁高皇帝遗朕平城之忧。高后时单于书绝悖逆。昔齐襄公复九世之雠。春秋大之。﹂是岁太初四年也。\\
\hspace*{2em}且鞮侯单于既立。尽归汉使之不降者。路充国等得归。单于初立。恐汉袭之。乃自谓﹁我兒子。安敢望汉天子。汉天子。我丈人行也﹂。汉遣中郎将苏武厚币赂遗单于。单于益骄。礼甚倨。非汉所望也。其明年。浞野侯破奴得亡归汉。\\
\hspace*{2em}其明年。汉使贰师将军广利以三万骑出酒泉。击右贤王於天山。得胡首虏万余级而还。匈奴大围贰师将军。几不脱。汉兵物故什六七。汉复使因杅将军敖出西河。与彊弩都尉会涿涂山。毋所得。又使骑都尉李陵将步骑五千人。出居延北千余里。与单于会。合战。陵所杀伤万余人。兵及食尽。欲解归。匈奴围陵。陵降匈奴。其兵遂没。得还者四百人。单于乃贵陵。以其女妻之。\\
\hspace*{2em}后二岁。复使贰师将军将六万骑。步兵十万。出朔方。彊弩都尉路博德将万余人。与贰师会。游击将军说将步骑三万人。出五原。因杅将军敖将万骑步兵三万人。出雁门。匈奴闻。悉远其累重於余吾水北。而单于以十万骑待水南。与贰师将军接战。贰师乃解而引归。与单于连战十余日。贰师闻其家以巫蛊族灭。因并众降匈奴。得来还千人一两人耳。游击说无所得。因杅敖与左贤王战。不利。引归。是岁汉兵之出击匈奴者不得言功多少。功不得御。有诏捕太医令随但。言贰师将军家室族灭。使广利得降匈奴。\\
\hspace*{2em}太史公曰。孔氏著春秋。隐桓之间则章。至定哀之际则微。为其切当世之文而罔襃。忌讳之辞也。世俗之言匈奴者。患其徼一时之权。而务谩纳其说。以便偏指。不参彼己。将率席中国广大。气奋。人主因以决策。是以建功不深。尧虽贤。兴事业不成。得禹而九州宁。且欲兴圣统。唯在择任将相哉。唯在择任将相哉。\\

% ===== 索隐述赞（司马贞《史记索隐》）· 统一样式 =====
% 题头「索隐述赞曰」五字并入反白框；正文每四句一列，定宽盒 \hfil 居列中；
% \fontsize{24}{41} 列距对齐全书界行（121+41k）；\clearpage 使述赞独占一页。
\clearpage
\shuzanhead{索隐述赞曰}\linebreak
{\shuzanbody\fontsize{24}{41}\selectfont
\hbox to \hsize{\hfil 獫狁。薰粥。居于北边。既称夏裔。\hfil}\linebreak
\hbox to \hsize{\hfil 式憬周篇。颇随畜牧。屡扰尘烟。爰自冒顿。\hfil}\linebreak
\hbox to \hsize{\hfil 尤聚控弦。虽空帑藏。未尽中权。\hfil}%
\par}
\end{document}
